\documentclass[11pt, a4paper]{article}

\usepackage[T1]{fontenc}
\usepackage[utf8]{inputenc}
\usepackage[a4paper, top=2.5cm, bottom=2.5cm, left=2.5cm, right=2.5cm]{geometry}
\usepackage{booktabs}
\usepackage{tabularx}
\usepackage{longtable}
\usepackage{enumitem}
\usepackage[hidelinks]{hyperref}
\usepackage{parskip}

\setlength{\parindent}{0pt}
\setlength{\parskip}{4pt plus 2pt}
\setlist[itemize]{topsep=2pt, itemsep=1pt, leftmargin=16pt}

\newcolumntype{L}[1]{>{\raggedright\arraybackslash}p{#1}}
\newcolumntype{C}[1]{>{\centering\arraybackslash}p{#1}}

\newcommand{\sechead}[1]{%
  \vspace{10pt}%
  \noindent{\large\bfseries #1}\\[-6pt]%
  \rule{\linewidth}{0.5pt}%
  \vspace{2pt}%
}

\pagestyle{plain}

\begin{document}

% ─── ENCABEZADO ────────────────────────────────────────────────────────
\begin{center}
{\large\bfseries PONTIFICIA UNIVERSIDAD JAVERIANA}\\[3pt]
Facultad de Ciencias Económicas y Administrativas\\[2pt]
\textit{Programa de Asignatura}\\[6pt]
{\LARGE\bfseries Estadística}\\[3pt]
% {\large Economía y Negocios Internacionales}
\end{center}

\vspace{6pt}
\rule{\linewidth}{0.5pt}
\vspace{2pt}

% ─── INFORMACIÓN GENERAL ───────────────────────────────────────────────
\sechead{Información General}

\begin{tabularx}{\linewidth}{L{4.5cm} X}
% \toprule
Departamento           & Ciencias Económicas y Administrativas \\
Período académico      & 2026-II \\
Intensidad horaria     & 3 horas semanales \\
Horario                & Viernes 7:00 -- 10:00 a.m. \\
Fechas del curso       & 31 de julio -- 27 de noviembre de 2026 \\
Sin clase              & 7 agosto (Batalla de Boyacá) y 28 agosto (Semana de Reflexión) \\
Profesor               & {Jaime Polanco Jimenez} \\
Correo                 & {jaime.polanco@javeriana.edu.co} \\
Campus virtual         & \url{https://campusvirtuallms.javeriana.edu.co} \\
Software               & R y RStudio \\
Datos                  & ICFES Saber 11 y Saber Pro (\url{icfes.gov.co}) \\
Repositorio            & \url{https://github.com/polanco-jaime/Curso-Estadistica-2026-3} \\
% \bottomrule
\end{tabularx}

% ─── OBJETIVOS ─────────────────────────────────────────────────────────
\sechead{Objetivos de Formación}

Al finalizar el curso el estudiante estará en capacidad de:

\begin{itemize}
  \item Aplicar técnicas de estadística descriptiva e inferencial sobre microdatos del ICFES para producir conclusiones relevantes.
  \item Comprender y usar las distribuciones muestrales fundamentales ($\chi^2$, $t$-Student y $F$-Fisher) y el papel que cumplen en la inferencia estadística.
  \item Estimar parámetros por máxima verosimilitud y mínimos cuadrados, e interpretar intervalos de confianza, errores estándar y diagnósticos de regresión.
  \item Especificar, estimar e interpretar modelos de regresión lineal simple, múltiple y logística.
  \item Diseñar, ejecutar e interpretar pruebas de hipótesis paramétricas y no paramétricas, incluyendo pruebas de bondad de ajuste.
  \item Desarrollar un proyecto integrador de análisis estadístico sobre datos ICFES, implementado en R.
\end{itemize}

% ─── CONTENIDOS MÍNIMOS ────────────────────────────────────────────────
\sechead{Contenidos Mínimos}

\begin{tabularx}{\linewidth}{L{3.5cm} X}
\toprule
\textbf{Unidad} & \textbf{Temas} \\
\midrule
U1 --- Estadística Descriptiva &
  Tendencia central y dispersión. Estadística gráfica. Muestreo y distribuciones muestrales.
  Distribuciones $\chi^2$, $t$-Student y $F$-Fisher. Correlación, EDA, Paradoja de Simpson. \\[4pt]
U2 --- Estimación Puntual &
  Propiedades de estimadores (insesgamiento, consistencia, eficiencia, MSE).
  Máxima Verosimilitud (MLE). Regresión lineal simple y múltiple (OLS).
  Diagnósticos de regresión. Regresión logística. \\[4pt]
U3 --- Estimación por Intervalo &
  IC para medias, proporciones y varianzas ($z$, $t$, $\chi^2$).
  IC bootstrap (percentil y BCa). Tamaño de muestra óptimo. \\[4pt]
U4 --- Pruebas de Hipótesis &
  Error tipo I y II, potencia estadística, $d$ de Cohen.
  Pruebas $t$, ANOVA, Chow. $\chi^2$ de independencia y bondad de ajuste.
  Pruebas no paramétricas: Mann-Whitney, Kruskal-Wallis. \\
\bottomrule
\end{tabularx}

% ─── METODOLOGÍA ───────────────────────────────────────────────────────
\sechead{Estrategias Metodológicas}

Cada sesión de viernes (7:00--10:00 a.m.) se divide en tres bloques de trabajo:

\begin{itemize}
  \item \textbf{Bloque 1 (7:00--8:00):} Exposición teórica con derivaciones y ejemplos.
  \item \textbf{Bloque 2 (8:10--9:10):} Aplicación práctica en R sobre microdatos ICFES . La variable central del curso
        es \texttt{MOD\_INGLES\_PUNT} como proxy de competitividad internacional.
  \item \textbf{Bloque 3 (9:20--10:00):} Taller grupal, debate o repaso.
        Retroalimentación de Problem Sets y del Proyecto Estadístico.
\end{itemize}

% ─── EVALUACIÓN ────────────────────────────────────────────────────────
\sechead{Evaluación}

\begin{tabularx}{\linewidth}{L{3.5cm} C{1.5cm} L{2.8cm} X}
\toprule
\textbf{Componente} & \textbf{Peso} & \textbf{Fecha} & \textbf{Descripción} \\
\midrule
Primer Parcial   & 25\% & Vie 4 sep 2026  &
  Unidad 1: descriptiva, muestreo, distribuciones $\chi^2$/$t$/$F$, correlación, EDA. \\[3pt]
Segundo Parcial  & 25\% & Vie 9 oct 2026  &
  Unidad 2: MLE, OLS simple y múltiple, diagnósticos, regresión logística. \\[3pt]
Examen Final     & 30\% & Vie 27 nov 2026 &
  Integración de las 4 unidades. Incluye ensayo de recomendación de política basada en
  datos Saber Pro NI. \\[3pt]
Proyecto Estadístico & 20\% & 3 entregas + 27 nov &
  Proyecto grupal (3--4 personas) en R con datos ICFES. Tres fases progresivas y
  presentación oral final. \\[3pt]
Problem Sets $\times 8$ & $+$0.5 bono & Semanal &
  Ocho talleres individuales sobre datos ICFES en R. Bono acumulativo de $+$0.5
  sobre la nota final. \\
\bottomrule
\end{tabularx}

Los exámenes permiten calculadora científica y una hoja A4 de fórmulas (ambas caras).

\medskip
\textit{Nota sobre el calendario:} el 28 de agosto (Semana de Reflexión) no hay clase.
Esa semana es ideal para avanzar en el PS3 y prepararse para el Primer Parcial del 4 de septiembre.

\newpage

% ─── PLAN TEMÁTICO ─────────────────────────────────────────────────────
\sechead{Plan Temático por Sesión}

\textit{Viernes 7:00--10:00 a.m. \quad 31 julio -- 27 noviembre 2026
\quad $|$ \quad Sin clase: 7 ago (festivo) y 28 ago (Semana de Reflexión)}

\vspace{6pt}
{\small
\setlength{\tabcolsep}{5pt}
\begin{longtable}{C{0.7cm} L{2.1cm} L{6.4cm} L{5.2cm}}
\toprule
\textbf{\#} & \textbf{Fecha} & \textbf{Tema} & \textbf{Actividades / Entregas} \\
\midrule
\endfirsthead
\multicolumn{4}{l}{\small\itshape (continuación)} \\
\toprule
\textbf{\#} & \textbf{Fecha} & \textbf{Tema} & \textbf{Actividades / Entregas} \\
\midrule
\endhead
\bottomrule
\endfoot
%
1  & Vie 31 jul &
  \textbf{U1 --- Descriptiva numérica.}
  Intro al curso y datos ICFES. Medidas de tendencia central y dispersión.
  Asimetría y curtosis con \texttt{MOD\_INGLES\_PUNT} en R. &
  Asigna PS1. Asigna Proyecto~F1. \\[3pt]
%
-- & \textit{Vie 7 ago} &
  \textit{Festivo --- Batalla de Boyacá. Sin clase.} & \\[3pt]
%
2  & Vie 14 ago &
  \textbf{U1 --- Descriptiva gráfica.}
  Histograma, KDE, box plot, scatter y heatmap departamental en R.
  Comparación IES pública vs.\ privada. &
  Entrega PS1. Asigna PS2. \\[3pt]
%
3  & Vie 21 ago &
  \textbf{U1 --- Muestreo, distribuciones y correlación.}
  TCL simulado con Saber 11. Las distribuciones $\chi^2(k)$, $t(k)$ y $F(k_1,k_2)$:
  de dónde vienen y cuándo aparecen. IC con $t$ vs.\ $z$ en R.
  Correlación Pearson y Spearman. Paradoja de Simpson. Pipeline EDA completo. &
  Entrega PS2. Asigna PS3. \\[3pt]
%
-- & \textit{Vie 28 ago} &
  \textit{Semana de Reflexión --- Sin clase.}
  \textit{Tiempo para PS3 y repaso antes del Parcial 1.} & \\[3pt]
%
E1 & Vie 4 sep &
  \textbf{Primer Parcial --- 25\%.}
  Parte I: preguntas conceptuales (MC + V/F + gráficos impresos).
  Parte II: dataset impreso ($n = 30$), cálculos y correlación.
  Parte III: interpretación de output R y diseño muestral. & \\[3pt]
%
4  & Vie 11 sep &
  \textbf{U2 --- Propiedades de estimadores y MLE.}
  Retroalimentación del Parcial~1. Insesgamiento, consistencia, eficiencia, MSE
  y bias-variance. Log-verosimilitud y \texttt{MASS::fitdistr()} en R. &
  Entrega PS3. Asigna PS4.
  \textbf{Entrega Proyecto~F1} (EDA, presentación oral 3~min). \\[3pt]
%
5  & Vie 18 sep &
  \textbf{U2 --- Regresión lineal simple (OLS).}
  Supuestos Gauss-Markov, derivación de $\hat{\beta}$, $R^2$, prueba $t$ y $F$.
  OLS en R: razonamiento $\sim$ inglés NI. Modelo con dummy \texttt{PRIVADA}. &
  Entrega PS4. Asigna PS5. \\[3pt]
%
6  & Vie 25 sep &
  \textbf{U2 --- Regresión lineal múltiple.}
  Sesgo por variable omitida (OVB), VIF, multicolinealidad, AIC/BIC.
  Modelo completo del inglés en NI en R. Interacción privada~$\times$~estrato. &
  Entrega PS5. Asigna PS6. \\[3pt]
%
7  & Vie 2 oct &
  \textbf{U2 --- Diagnósticos y regresión logística.}
  Breusch-Pagan, Durbin-Watson, Cook's $D$, errores estándar robustos.
  Logit para \texttt{supera\_B2}, odds ratios y AUC-ROC en R. &
  Entrega PS6.
  \textbf{Entrega Proyecto~F2} (modelo OLS/logit + IC, 3~min oral). \\[3pt]
%
E2 & Vie 9 oct &
  \textbf{Segundo Parcial --- 25\%.}
  Parte I: estimadores y supuestos OLS.
  Parte II: regresión múltiple con datos ICFES NI (coeficientes, VIF, OVB).
  Parte III: diagnósticos + logit + ensayo. & \\[3pt]
%
8  & Vie 16 oct &
  \textbf{U3 --- Intervalos de confianza.}
  Retroalimentación del Parcial~2. Lógica frecuentista, IC con $z$, $t$ y $\chi^2$.
  IC para proporciones. Forest plot de 33 departamentos en R. &
  Asigna PS7 (2 semanas). \\[3pt]
%
9  & Vie 23 oct &
  \textbf{U3 --- Bootstrap y tamaño de muestra.}
  IC bootstrap percentil y BCa con \texttt{boot::boot()} en R.
  Fórmula de $n$ óptimo. Muestreo estratificado y costo-eficiencia. & \\[3pt]
%
10 & Vie 30 oct &
  \textbf{U4 --- Errores y potencia estadística.}
  $H_0$ y $H_1$, $p$-valor (errores frecuentes de interpretación).
  Error tipo I y II, potencia $= 1 - \beta$, $d$ de Cohen. Curva de potencia en R. &
  Entrega PS7. Asigna PS8 (2 semanas). \\[3pt]
%
11 & Vie 6 nov &
  \textbf{U4 --- Pruebas paramétricas.}
  Prueba $t$ (1 muestra, 2 muestras, pareada). ANOVA + Tukey.
  Prueba de Chow: quiebre estructural en el inglés de NI antes y después de la pandemia. & \\[3pt]
%
12 & Vie 13 nov &
  \textbf{U4 --- Pruebas no paramétricas y bondad de ajuste.}
  $\chi^2$ de independencia y bondad de ajuste (¿es normal \texttt{MOD\_INGLES\_PUNT}?).
  Mann-Whitney, Kruskal-Wallis. Tablas de contingencia con dos criterios en R. &
  Entrega PS8. \\[3pt]
%
13 & Vie 20 nov &
  \textbf{Repaso final e instrucciones.}
  Práctica de problemas integradores (tipo examen) de las 4 unidades.
  Retroalimentación colectiva de PS8.
  Ensayo de presentaciones del proyecto (3~min/grupo). &
  Ensayo de presentaciones. Logística del examen final. \\[3pt]
%
F  & Vie 27 nov &
  \textbf{Sesión de cierre --- Examen Final (30\%) + Presentaciones (20\%).}
  7:00--8:20: presentaciones del proyecto (10~min/grupo + preguntas).
  8:20--8:30: descanso.
  8:30--10:00: examen final integrado (4 unidades). &
  Examen Final 30\%.
  \textbf{Entrega Proyecto~F3: informe escrito ($\leq$15~págs.) al inicio.} \\
%
\end{longtable}
}

\newpage

% ─── PROBLEM SETS ──────────────────────────────────────────────────────
\sechead{Problem Sets --- Trabajo Autónomo}

Los Problem Sets son individuales y se resuelven en R sobre microdatos ICFES. Cada
entrega completa suma $\mathbf{+0.0625}$ a la nota final (total si entrega los 8: $\mathbf{+0.5}$).
Entrega digital al inicio de la clase indicada. El código R debe incluirse.

\vspace{6pt}
{\small
\begin{tabularx}{\linewidth}{C{0.8cm} L{1.8cm} L{1.8cm} X}
\toprule
\textbf{PS} & \textbf{Asignado} & \textbf{Entrega} & \textbf{Enunciado} \\
\midrule
1 & Jul 31 (Ses 1) & Ago 14 (Ses 2) &
  Análisis descriptivo completo de \texttt{MOD\_INGLES\_PUNT} en NI: media, mediana, DE,
  asimetría, curtosis y cuartiles por departamento, tipo de IES y estrato. Mínimo 5
  visualizaciones en R. \\[4pt]
2 & Ago 14 (Ses 2) & Ago 21 (Ses 3) &
  Simulación del TCL con 2000 muestras de \texttt{PUNT\_INGLES} del Saber 11 para
  $n = 30, 100, 500$. Graficar convergencia a $N(\mu,\,\sigma/\sqrt{n})$
  y calcular el error estándar teórico vs.\ simulado. \\[4pt]
3 & Ago 21 (Ses 3) & Sep 11 (Ses 4) &
  Matriz de correlaciones Pearson y Spearman de los 5 módulos Saber Pro en NI.
  Tabla de contingencia género $\times$ nivel inglés MCER.
  Demostración de la Paradoja de Simpson con datos ICFES.
  \textit{Plazo extendido: la Semana de Reflexión (28~ago) y el Parcial~1 (4~sep)
  permiten tres semanas para este PS.} \\[4pt]
4 & Sep 11 (Ses 4) & Sep 18 (Ses 5) &
  MLE en R (\texttt{MASS::fitdistr}) para distribución de \texttt{MOD\_RAZONA\_CUANTITAT\_PUNT}
  en NI: comparar normal vs.\ $t$ vs.\ beta por AIC. OLS simple: ¿predice el inglés
  del Saber 11 el razonamiento cuantitativo del Saber Pro? \\[4pt]
5 & Sep 18 (Ses 5) & Sep 25 (Ses 6) &
  OLS múltiple del inglés en NI: \texttt{ESTRATO + PRIVADA + log(INGRESOS) + HORAS + REGION}.
  Interpretar cada $\hat{\beta}$, calcular VIF y mostrar el sesgo por variable omitida
  comparando el modelo simple con el múltiple. \\[4pt]
6 & Sep 25 (Ses 6) & Oct 2 (Ses 7) &
  Diagnósticos del modelo OLS en R: 4 gráficos de residuos, Breusch-Pagan,
  Durbin-Watson, Cook's $D$. Identificar los 3 puntos más influyentes y aplicar
  errores estándar robustos (Huber-White). \\[4pt]
7 & Oct 16 (Ses 8) & Oct 30 (Ses 10) &
  IC bootstrap BCa ($B = 2000$) al 95\% para la brecha pública--privada en inglés
  NI por macro-región. ¿Es significativa en todas las regiones? Calcular el $n$
  óptimo para $E = \pm 2$ puntos al 95\%. \\[4pt]
8 & Oct 30 (Ses 10) & Nov 13 (Ses 12) &
  (1) $\chi^2$ de independencia: nivel inglés MCER $\times$ tipo IES, con $V$ de Cramér.
  (2) Mann-Whitney $+$ $d$ de Cohen: ¿difiere el inglés por género en NI?
  (3) Chow o Kruskal-Wallis: ¿hubo quiebre en el inglés en NI con la pandemia? \\
\bottomrule
\end{tabularx}
}

% ─── PROYECTO ESTADÍSTICO ──────────────────────────────────────────────
\sechead{Proyecto Integrador de Análisis Estadístico (20\%)}

Grupos de 3--4 personas, conformados en la primera sesión. El proyecto se desarrolla
íntegramente en R con microdatos ICFES. El informe escrito final ($\leq 15$~páginas)
se entrega el 27 de noviembre al inicio de la clase.

\vspace{6pt}
\begin{tabularx}{\linewidth}{L{3.2cm} C{1.5cm} L{2.8cm} X}
\toprule
\textbf{Fase} & \textbf{Peso} & \textbf{Entrega} & \textbf{Contenido} \\
\midrule
Fase 1 --- EDA & 20\% & \textbf{Sep 11} (Ses 4)\newline oral 3 min &
  Limpieza y exploración del dataset ICFES elegido: estadísticos descriptivos,
  $\geq 6$ visualizaciones (\texttt{ggplot2}), tratamiento de outliers y 3 hallazgos
  clave. La Semana de Reflexión (28~ago) es tiempo útil para esta fase. \\[4pt]
Fase 2 --- Modelo + IC & 35\% & \textbf{Oct 2} (Ses 7)\newline oral 3 min &
  OLS múltiple o regresión logística en R, diagnósticos completos (residuos, VIF,
  Breusch-Pagan), IC al 95\% para coeficientes clave. Evaluación: $R^2$ o AUC. \\[4pt]
Fase 3 --- Pruebas + Informe & 45\% & \textbf{Nov 27} (Ses F)\newline oral 10 min &
  Pruebas de hipótesis sobre los hallazgos ($t$, ANOVA, $\chi^2$ o Mann-Whitney),
  análisis de potencia, $d$ de Cohen, implicaciones para NI.
  Presentación oral 10~min $+$ informe escrito entregado al inicio de la sesión. \\
\bottomrule
\end{tabularx}

\vspace{6pt}
\textbf{Preguntas de investigación sugeridas:}
(1)~¿Qué determina el nivel de inglés al graduarse en NI?
(2)~¿Qué variables predicen superar el B2?
(3)~¿Ha cerrado la brecha de género en inglés en NI entre 2014 y 2024?
(4)~¿Cuánto valor agrega la universidad sobre el perfil de entrada del estudiante?
(5)~¿Difieren significativamente los promedios de inglés entre macro-regiones?

% ─── REGLAS ────────────────────────────────────────────────────────────
\sechead{Reglas del Curso}

\begin{itemize}
  \item \textbf{Asistencia.} La asistencia es fundamental. Cada sesión integra teoría,
        aplicación en R y taller; su contenido no se recupera con notas de compañeros.
  \item \textbf{Problem Sets.} Son individuales. Entrega digital al inicio de la clase
        indicada. Las entregas tardías pierden el bono. El código R debe incluirse.
  \item \textbf{Proyecto.} Los grupos no se modifican una vez conformados.
        Cada fase es condición para la siguiente. El código R se entrega como anexo.
  \item \textbf{Reclamos.} En la sesión en que se devuelve la evaluación, con argumento
        claro. Los reclamos de PS se hacen con el monitor en su horario.
  \item \textbf{Honestidad académica.} Cualquier copia o intento de copia se sancionará
        según el Reglamento Estudiantil (Párrs.~118, 123-d y~128):
        pérdida de la asignatura con nota~0.0.
  \item \textbf{Exámenes.} Calculadora científica y una hoja A4 de fórmulas (ambas caras).
        Sin computadores ni celulares. Exámenes individuales.
\end{itemize}

% ─── BIBLIOGRAFÍA ──────────────────────────────────────────────────────
\sechead{Fuentes de Información}

\textbf{Textos guía}
\begin{itemize}
  \item Canavos, G. \textit{Probabilidad y estadística, aplicación y métodos}. McGraw-Hill, 1988.
  \item Freund, J.; Miller, I.; Miller, M. \textit{Estadística Matemática con Aplicaciones}.
        6.ª ed. Prentice Hall.
  \item Wackerly, D.; Mendenhall, W.; Sheaffer, R. \textit{Estadística Matemática con
        Aplicaciones}. Thomson.
\end{itemize}

\begin{sloppypar}
\textbf{Software y datos (gratuitos)}
\begin{itemize}
  \item R Project: \url{https://www.r-project.org} \quad RStudio: \url{https://posit.co}
  \item Paquetes clave: \texttt{tidyverse}, \texttt{ggplot2}, \texttt{car},
        \texttt{lmtest}, \texttt{sandwich}, \texttt{MASS}, \texttt{boot},
        \texttt{pwr}, \texttt{effsize}, \texttt{strucchange}
  \item ICFES microdatos: \url{https://icfes.gov.co/resultados-saber}
  \item Canal YouTube del curso anterior (Prof.\ Ronderos):
        \url{https://www.youtube.com/@nicolasronderos615}
\end{itemize}
\end{sloppypar}

\textbf{Referencias avanzadas}
\begin{itemize}
  \item Casella, G.; Berger, R. \textit{Statistical Inference}. 2.ª ed. Duxbury Press, 2002.
  \item DeGroot, M.\ H. \textit{Probability and Statistics}. 2.ª ed. Addison-Wesley, 1989.
\end{itemize}

\vspace{16pt}
\begin{center}
\rule{0.5\linewidth}{0.4pt}\\[4pt]
{\small Pontificia Universidad Javeriana $\cdot$
  Facultad de Ciencias Económicas y Administrativas $\cdot$ 2026-II}
\end{center}

\end{document}